```latex
\documentclass[12pt]{article}
\usepackage[utf8]{inputenc}
\usepackage{amsmath, amssymb, graphicx, hyperref}
\usepackage{geometry}
\geometry{a4paper, margin=1in}
\title{SHOA-Adaptive: Programmable Geopolymer Self-Assembly via GRAPE-Optimized Acoustic Resonance with Formal Validation Criteria}
\author{Konstantin Fedotov}
\date{\today}

\begin{document}

\maketitle

\begin{abstract}
We present SHOA-Adaptive, a new class of materials -- geopolymers capable of controlled self-assembly under a GRAPE-optimized ultrasonic pulse. The key innovation is the introduction of SHOA centers (superparamagnetic nanoparticles) into a mineral matrix, functioning as quantum sensors and actuators governed by the Spiked Voting protocol. We provide formal, testable criteria for experimental verification, including tests for assembly precision, energy efficiency, cycle stability, and the distinction between quantum-controlled assembly and conventional ultrasonic compaction.
\end{abstract}

\section{Introduction}
The SHOA project has demonstrated self-healing (Classic), learning (Neuro), and neural integration (NeuroSHOA). SHOA-Adaptive introduces programmable matter: geopolymers that self-assemble under quantum control. This work provides the formal experimental criteria needed to verify that the observed assembly is genuinely controlled by the GRAPE-optimized pulse, and not by conventional acoustic effects.

\section{Model}
\subsection{Full Hamiltonian}
The total Hamiltonian of an ensemble of SHOA centers embedded in a viscoelastic matrix is:
\begin{equation}
\mathcal{H}_{\text{Adaptive}} = \mathcal{H}_{\text{spin}} + \mathcal{H}_{\text{dip}} + \mathcal{H}_{\text{acoustic}}(t) + \mathcal{H}_{\text{geometry}},
\end{equation}
where the key extension is $\mathcal{H}_{\text{acoustic}}(t) = \sum_i \Omega_i(t) \sigma_z^{(i)}$, a GRAPE-optimized ultrasonic pulse that modulates the energy gap of each center. The dynamics of the matrix are described by a self-consistent equation linking quantum force and classical elastic relaxation.

\subsection{GRAPE Optimization}
The pulse shape $\Omega_i(t)$ is optimized using gradient ascent to maximize the fidelity of the target geometry. The optimization loop runs over 5000 iterations with convergence tolerance $10^{-6}$.

\section{Formal Validation Criteria}

\subsection{Assembly Precision (Criterion I)}
\begin{itemize}
    \item \textbf{Definition:} The final shape of the geopolymer block must match the target shape with a deviation of less than 5\% of the characteristic dimension.
    \item \textbf{What to measure:} The root-mean-square deviation (RMSD) between the scanned surface profile of the cured block and the digital reference model.
    \item \textbf{Target:} RMSD $< 5\%$ of the block's linear dimension for a GRAPE-optimized pulse; RMSD $> 15\%$ for an unoptimized monochromatic pulse of the same energy.
    \item \textbf{Instrumentation:} Laser profilometry or structured-light 3D scanning.
\end{itemize}

\subsection{Energy Efficiency (Criterion II)}
\begin{itemize}
    \item \textbf{Definition:} The GRAPE-optimized pulse must achieve the target shape using less than 50\% of the acoustic energy required by a conventional ultrasonic compaction.
    \item \textbf{What to measure:} The total acoustic energy delivered to the sample during the assembly process.
    \item \textbf{Target:} $E_{\text{GRAPE}} < 0.5 \times E_{\text{conventional}}$ for the same final RMSD.
    \item \textbf{Significance:} This criterion proves that the GRAPE optimization provides genuine efficiency gains through resonance, not just brute-force compaction.
\end{itemize}

\subsection{Cycle Stability and Learning (Criterion III)}
\begin{itemize}
    \item \textbf{Definition:} The assembly precision must improve over repeated cycles as the material ``learns'' its optimal response.
    \item \textbf{What to measure:} RMSD over 10 consecutive assembly-disassembly cycles.
    \item \textbf{Target:} RMSD must decrease monotonically and saturate, with a learning rate $\lambda > 0.01$ when fitted to the memory law $K_{\text{eff}}(N) = K_0 (1 - \beta e^{-\lambda N})$.
    \item \textbf{Control:} A sample without SHOA centers must show no significant improvement over cycles.
\end{itemize}

\subsection{Distinguishability from Conventional Compaction (Criterion IV)}
\begin{itemize}
    \item \textbf{Definition:} The GRAPE-optimized assembly must produce a different final microstructure than conventional ultrasonic compaction.
    \item \textbf{What to measure:} Scanning electron microscopy (SEM) images of the internal microstructure: porosity, grain alignment, and nanoparticle distribution.
    \item \textbf{Target:} The GRAPE sample must show ordered alignment of nanoparticles along the acoustic axis, while the conventional sample shows random distribution.
    \item \textbf{Significance:} This criterion proves that SHOA-Adaptive is not just a rebranding of existing techniques, but a fundamentally new assembly mechanism.
\end{itemize}

\section{Experimental Protocol}
We propose a concrete, reproducible experiment:
\begin{enumerate}
    \item \textbf{Synthesis:} Superparamagnetic Fe$_3$O$_4$ nanoparticles ($\sim$10-15 nm) are introduced into a geopolymer mixture (metakaolin + alkaline solution).
    \item \textbf{Scanning:} The geometry of a reference block is read by a laser profilometer.
    \item \textbf{Activation:} A GRAPE-optimized ultrasonic pulse is applied to the mixture, causing adaptive flow.
    \item \textbf{Fixation:} After curing, the block's shape is compared to the reference.
\end{enumerate}

\section{Conclusion}
SHOA-Adaptive provides a formally verifiable platform for programmable matter. The four criteria presented here transform the concept of quantum-controlled geopolymer assembly into a testable engineering discipline.

\end{document}
